\section{Conclusion}
\label{sec:conclusion}

Iterating AI \emph{review $\to$ revise} does \emph{not} reach a fixed point. The regenerated
critique injects endless edit pressure; the document inflates $\sim 7\times$, drifts from its
original meaning (cosine $0.97\to 0.77$), and settles into a high, non-zero churn plateau
(mean final $\Delta\approx 0.38$) rather than a stable point, with $0\%$ of trajectories
converging. Removing the separate review stage restores convergence to a soft fixed point
($42\%$ approximate fixed points, $R^2=0.98$), identifying the explicit review loop---not
iteration itself---as the cause of non-convergence. Along the way we document a concrete
failure of isolated LLM-judges: quality scores rise while meaning preservation collapses.

\para{Takeaway.} ``Keep asking an AI to review and revise until it stabilizes'' is not a
reliable stopping strategy; if a fixed point is desired, either drop the separate review stage
or add an explicit convergence/stop criterion.

\para{Future work.} We highlight five directions: (1) a \emph{positive control for
convergence}---give the reviewer an explicit ``return \textsc{no changes needed}'' option
and/or the reviser a hard length budget, and test whether a fixed point reappears;
(2) \emph{cross-provider replication} (Claude, Gemini, open-weight) to establish generality of
the never-satisfied-reviewer mechanism; (3) a \emph{length-controlled operator} to disentangle
churn from genuine rewriting versus churn from growth; (4) \emph{basins of attraction}---re-seed
the same content as paraphrases and test whether trajectories share an attractor; and
(5) \emph{human evaluation} of meaning drift to validate the LLM-judge fidelity signal.
